\documentclass[aspectratio=169]{beamer}
\usetheme{metropolis}
\usepackage{nexus_theme}

\title{Predatory Publishing}
\subtitle{Module 0.4: Spotting the Scams of the Academic World}
\author{Nexus Scholar Suite --- Modern Research Methodologies}
\date{}

\begin{document}

% --- TITLE SLIDE ---
\begin{frame}[plain]
  \titlepage
\end{frame}

% --- LEARNING OBJECTIVES ---
\begin{frame}{What You'll Learn}
  By the end of this lesson, you will be able to:
  \begin{enumerate}
    \item \textbf{Define} predatory publishing and its financial incentives.
    \item \textbf{Identify} the common red flags of a predatory journal.
    \item \textbf{Explain} the difference between legitimate Open Access APCs and predatory fees.
    \item \textbf{Utilize} standard whitelists (DOAJ) to verify venues.
  \end{enumerate}
\end{frame}

% --- THE SCAM ---
\begin{frame}{The Business Model of Deception}
  \begin{alertblock}{How it Works}
    \begin{itemize}
      \item \textbf{The Exploitation:} Researchers are desperate to publish to keep their jobs ("Publish or Perish").
      \item \textbf{The Scam:} The publisher charges an Article Processing Charge (APC) of \$500-\$2,000.
      \item \textbf{The Catch:} They provide \textbf{zero} peer review and \textbf{zero} editorial oversight. They will publish literally anything.
    \end{itemize}
  \end{alertblock}
  \vspace{0.3cm}
  \textit{Consequence:} Publishing here permanently stains a junior researcher's CV and pollutes the scientific record.
\end{frame}

% --- LEGITIMATE VS PREDATORY ---
\begin{frame}{The Open Access Confusion}
  \begin{columns}[T]
    \begin{column}{0.48\textwidth}
      \begin{block}{Legitimate Open Access (OA)}
        \begin{itemize}
          \item e.g., \textit{PLOS ONE}, \textit{Nature Comms}
          \item Charges an APC fee.
          \item Rigorous, months-long peer review.
          \item High rejection rate.
          \item Science is free for the public to read.
        \end{itemize}
      \end{block}
    \end{column}
    \begin{column}{0.48\textwidth}
      \begin{alertblock}{Predatory Open Access}
        \begin{itemize}
          \item e.g., \textit{Int'l Journal of Fake Science}
          \item Charges an APC fee.
          \item \textbf{NO} peer review.
          \item 100\% acceptance rate.
          \item Exists solely to collect the fee.
        \end{itemize}
      \end{alertblock}
    \end{column}
  \end{columns}
\end{frame}

% --- RED FLAGS ---
\begin{frame}{The Red Flags of a Predatory Journal}
  \begin{exampleblock}{If you see these, run away:}
    \begin{enumerate}
      \item \textbf{Spam Emails:} "Dear Esteemed Dr. Smith, we urgently need your paper for our next issue!"
      \item \textbf{Rapid Acceptance:} "Guaranteed peer review and publication in 7 days." (Real review takes months).
      \item \textbf{Bizarre Scope:} "The Global Journal of Dentistry, Astrophysics, and Poetry."
      \item \textbf{Fake Metrics:} Claiming a "Global Impact Factor" (a fake metric designed to sound like the real Clarivate Impact Factor).
    \end{enumerate}
  \end{exampleblock}
\end{frame}

% --- TOOLS FOR DEFENSE ---
\begin{frame}{Defending Yourself}
  \begin{block}{Whitelists and Blacklists}
    \begin{itemize}
      \item \textbf{DOAJ (Directory of Open Access Journals):} The definitive \textit{whitelist}. If an OA journal is not listed here, be highly suspicious.
      \item \textbf{Think. Check. Submit.:} An industry-standard checklist for researchers to evaluate venues.
      \item \textbf{Beall's List:} The original \textit{blacklist}. Officially taken down due to legal threats, but anonymous archives exist online.
    \end{itemize}
  \end{block}
\end{frame}

% --- KEY TAKEAWAY ---
\begin{frame}{Key Takeaway}
  \begin{block}{What We Accomplished}
    \begin{enumerate}
      \item \textbf{The Scam}: Predatory journals sell fake peer review.
      \item \textbf{The Nuance}: Fees don't mean fake. Lack of review means fake.
      \item \textbf{The Defense}: Check the DOAJ before submitting.
    \end{enumerate}
  \end{block}
  \vspace{0.3cm}
  \textbf{Next:} Module 1.1 --- Reference Management 101
\end{frame}

% --- EXERCISE PREVIEW ---
\begin{frame}{Your Turn: Exercise}
  \begin{exampleblock}{Exercise: The Predatory Audit}
    Act as an academic detective. Audit three unsolicited emails from "journals" against our red flags and determine which belong in the trash.
  \end{exampleblock}
\end{frame}

% --- STANDOUT CLOSE ---
\begin{frame}[standout]
  Protecting your reputation is just as important\\
  as publishing your research.
\end{frame}

\end{document}
